% 函数类追随随机符号：三个函数在5个样本点上的取值与两串随机符号的内积。
% 数据：f1=(1,1,1,1,1), f2=(1,1,-1,-1,1), f3=(-1,1,1,-1,-1)；
% sigma^(a)=(+,-,-,+,+) 给出内积 (1,1,-5)，最大值1；
% sigma^(b)=(+,+,-,-,-) 给出内积 (-1,3,1)，最大值3。
% 对全部32串符号平均 max_j <sigma,f_j>/5 得 0.4（枚举复算）。
\begingroup
\providecolor{FigureGreen}{HTML}{829C8B}
\providecolor{FigureGreenFill}{HTML}{E8EEE9}
\begin{tikzpicture}[figure picture,x=1mm,y=1mm]
\def\cw{9}\def\ch{6}
% 列标题
\foreach \i in {1,...,5}{\node[figure note] at ({18+\i*\cw},47) {$z_{\i}$};}
\node[figure note] at (82,47) {$\langle\bm\sigma^{(a)},f\rangle$};
\node[figure note] at (101,47) {$\langle\bm\sigma^{(b)},f\rangle$};
% 函数行
\def\rowf#1#2#3{%
 \node[anchor=east] at (24,#1) {$#2$};
 \foreach \v [count=\i] in {#3}{%
  \ifnum\v>0
   \filldraw[draw=FigureBlue,fill=FigureBlueFill,line width=.5pt] ({18+\i*\cw-\cw/2+.6},{#1-\ch/2+.6}) rectangle ({18+\i*\cw+\cw/2-.6},{#1+\ch/2-.6});
   \node[text=FigureBlue] at ({18+\i*\cw},#1) {$+1$};
  \else
   \filldraw[draw=FigureRed,fill=FigureRedFill,line width=.5pt] ({18+\i*\cw-\cw/2+.6},{#1-\ch/2+.6}) rectangle ({18+\i*\cw+\cw/2-.6},{#1+\ch/2-.6});
   \node[text=FigureRed] at ({18+\i*\cw},#1) {$-1$};
  \fi}}
\rowf{40}{f_1}{1,1,1,1,1}
\rowf{33}{f_2}{1,1,-1,-1,1}
\rowf{26}{f_3}{-1,1,1,-1,-1}
% 内积列
\node at (82,40) {$1$};\node at (82,33) {$1$};\node at (82,26) {$-5$};
\node at (101,40) {$-1$};\node at (101,26) {$1$};
\draw[FigureGreen,line width=.9pt,rounded corners=1mm] (96.5,30.2) rectangle (105.5,35.8);
\node at (101,33) {$3$};
\draw[FigureGreen,line width=.9pt,rounded corners=1mm] (77.5,37.2) rectangle (86.5,42.8);
\draw[FigureGreen,line width=.9pt,rounded corners=1mm] (77.5,30.2) rectangle (86.5,35.8);
\draw[FigureGrid,line width=.5pt] (20,21.5)--(108,21.5);
% 符号行
\def\rows#1#2#3{%
 \node[anchor=east] at (24,#1) {$#2$};
 \foreach \v [count=\i] in {#3}{%
  \filldraw[draw=FigureYellow,fill=FigureYellowFill,line width=.5pt] ({18+\i*\cw-\cw/2+.6},{#1-\ch/2+.6}) rectangle ({18+\i*\cw+\cw/2-.6},{#1+\ch/2-.6});
  \node at ({18+\i*\cw},#1) {$\v$};}}
\rows{16}{\bm\sigma^{(a)}}{+,-,-,+,+}
\rows{9}{\bm\sigma^{(b)}}{+,+,-,-,-}
\node[figure note,anchor=west] at (72,16) {最大内积 $1$，贡献 $1/5$};
\node[figure note,anchor=west] at (72,9) {最大内积 $3$，贡献 $3/5$};
\node[figure note] at (64,1.5) {对全部 $2^5=32$ 串符号平均：$\widehat{\mathfrak R}_S=0.4$};
\end{tikzpicture}
\endgroup
